% Source: conditional_random_fields.tex
\begin{notechapter}
\chapter{Conditional Random Fields for Part-of-Speech Tagging}\label{ch:crf}

\section{Introduction}
Conditional Random Fields (CRFs) are a type of 
discriminative probabilistic model used for sequence 
labeling tasks, such as Part-of-Speech (POS) tagging, 
named entity recognition (NER), and more. Unlike Hidden 
Markov Models (HMMs), which model the joint probability 
\( P(x, y) \), CRFs directly model the conditional 
probability \( P(y | x) \), making them more flexible 
and powerful for many tasks.

\section{Key Components of CRFs}
\begin{itemize}
    \item Observations: A sequence of input features 
          (e.g., words in a sentence).
    \item Labels: A sequence of output labels 
          (e.g., POS tags).
    \item Feature Functions: Functions that capture 
          relationships between observations and labels.
    \item Parameters: Weights associated with feature 
          functions, learned during training.
\end{itemize}

\section{Mathematical Formulation}
A CRF defines the conditional probability of a label 
sequence \( \mathbf{y} = (y_1, y_2, \dots, y_T) \) given 
an observation sequence \( \mathbf{x} = (x_1, x_2, \dots, 
x_T) \) as:
\[
P(\mathbf{y} | \mathbf{x}) = \frac{1}{Z(\mathbf{x})} 
\exp\left( \sum_{t=1}^T \sum_{k=1}^K \lambda_k f_k(y_{t-1}, 
y_t, \mathbf{x}, t) \right),
\]
where:
\begin{itemize}
    \item \( f_k(y_{t-1}, y_t, \mathbf{x}, t) \): A 
          feature function that captures dependencies 
          between labels and observations.
    \item \( \lambda_k \): The weight associated with 
          the \( k \)-th feature function.
    \item \( Z(\mathbf{x}) \): The partition function, 
          ensuring the probabilities sum to 1.
\end{itemize}

\section{Feature Functions}
Feature functions in CRFs can capture a wide range of 
dependencies, such as:
\begin{itemize}
    \item Transition Features: Dependencies between 
          adjacent labels (e.g., \( y_{t-1} = \text{Det} \) 
          and \( y_t = \text{Noun} \)).
    \item Emission Features: Dependencies between a 
          label and the current observation (e.g., 
          \( y_t = \text{Noun} \) and \( x_t = \text{"cat"} \)).
    \item Contextual Features: Dependencies between a 
          label and the surrounding observations (e.g., 
          \( y_t = \text{Verb} \) and \( x_{t-1} = \text{"the"} \)).
\end{itemize}

\section{Graphical Representation}
CRFs are typically represented as an undirected graph, 
where nodes represent labels and edges represent 
dependencies between them. Below is an example of a 
linear-chain CRF for POS tagging:

\begin{center}
\begin{tikzpicture}[
    > = Stealth, % Arrow head style
    shorten > = 1pt, % Don't touch arrow head to node
    auto,
    node distance = 2cm, % Distance between nodes
    semithick, % Line style
]

    % Nodes for labels
    \node[state, color=blue] (y1) {\( y_1 \)};
    \node[state, color=blue] (y2) [right of=y1] {\( y_2 \)};
    \node[state, color=blue] (y3) [right of=y2] {\( y_3 \)};
    \node[state, color=blue] (y4) [right of=y3] {\( y_4 \)};

    % Edges for dependencies
    \path[-, color=blue] (y1) edge node {} (y2);
    \path[-, color=blue] (y2) edge node {} (y3);
    \path[-, color=blue] (y3) edge node {} (y4);

    % Nodes for observations
    \node[state, color=red] (x1) [below of=y1, yshift=-0.1cm] {\( x_1 \)};
    \node[state, color=red] (x2) [below of=y2, yshift=-0.1cm] {\( x_2 \)};
    \node[state, color=red] (x3) [below of=y3, yshift=-0.1cm] {\( x_3 \)};
    \node[state, color=red] (x4) [below of=y4, yshift=-0.1cm] {\( x_4 \)};

    % Edges for emission features
    \path[-, color=red] (y1) edge node {} (x1);
    \path[-, color=red] (y2) edge node {} (x2);
    \path[-, color=red] (y3) edge node {} (x3);
    \path[-, color=red] (y4) edge node {} (x4);
\end{tikzpicture}
\end{center}

\section{Comparison with HMMs}
\begin{itemize}
    \item Generative vs. Discriminative: HMMs are 
          generative models that model \( P(x, y) \), 
          while CRFs are discriminative models that 
          directly model \( P(y | x) \).
    \item Feature Flexibility: CRFs can incorporate 
          arbitrary features of the input sequence, 
          while HMMs are limited to emission and 
          transition probabilities.
    \item Training: CRFs are trained using maximum 
          likelihood estimation (MLE) or gradient-based 
          methods, while HMMs use the Baum-Welch 
          algorithm (a variant of EM).
\end{itemize}

\section{Example: POS Tagging with CRFs}
Consider the sentence: \textit{"The cat sat on the floor"}. 
The goal is to predict the POS tags for each word.

\subsection{Feature Functions}
\begin{itemize}
    \item Transition Features:
        \begin{itemize}
            \item \( f_1(y_{t-1} = \text{Det}, y_t = \text{Noun}) = 1 \)
            \item \( f_2(y_{t-1} = \text{Noun}, y_t = \text{Verb}) = 1 \)
        \end{itemize}
    \item Emission Features:
        \begin{itemize}
            \item \( f_3(y_t = \text{Det}, x_t = \text{"The"}) = 1 \)
            \item \( f_4(y_t = \text{Noun}, x_t = \text{"cat"}) = 1 \)
        \end{itemize}
\end{itemize}

\subsection{Inference}
Inference in CRFs involves finding the most likely 
label sequence \( \mathbf{y}^* \) given the 
observations \( \mathbf{x} \):
\[
\mathbf{y}^* = \arg\max_{\mathbf{y}} P(\mathbf{y} | \mathbf{x}).
\]
This is typically done using the Viterbi algorithm, 
adapted for CRFs.

\section{Conclusion}
Conditional Random Fields are a powerful and flexible 
model for sequence labeling tasks. By directly modeling 
\( P(y | x) \) and incorporating rich feature functions, 
CRFs often outperform generative models like HMMs in 
practice. Their ability to handle arbitrary dependencies 
makes them a popular choice for NLP tasks.
\end{notechapter}
